\section{Presentación de la empresa}

\subsection{Reseña de la trayectoria, capacidades instaladas y líneas de negocio}

\textbf{audIT} es una empresa de ingeniería de software, arquitectura de datos y ciberseguridad industrial fundada en 2022 (4 años de trayectoria comprobable en el mercado nacional), orientada exclusivamente a resolver la continuidad operacional, la visibilidad telemática y la trazabilidad analítica en organizaciones con operaciones geográficamente distribuidas y condiciones de conectividad intermitente. Su propuesta de valor se concentra en un segmento donde la mayoría de los proveedores generalistas falla: integrar dispositivos físicos desplegados en campo ---sensores, equipos embarcados, telemetría vehicular--- con plataformas de datos que deben operar de forma confiable incluso sin conexión permanente.

\begin{itemize}
    \item \textbf{Giro Principal:} Servicios integrales de consultoría informática, diseño e implantación de plataformas de datos en la nube, arquitectura de software híbrida, desarrollo de soluciones IoT / telemetría de terreno y auditoría de ciberseguridad bajo estándares internacionales.
    \item \textbf{Misión:} Diseñar e implantar soluciones de software e integración ciber-física de alta confiabilidad que transformen operaciones críticas de campo en flujos de datos auditables, continuos y seguros, habilitando la toma de decisiones basada en evidencia técnica rigurosa.
    \item \textbf{Visión:} Consolidarse como el socio tecnológico de referencia en el Cono Sur para la modernización de sistemas de misión crítica y la integración de cadenas logísticas y de transporte bajo arquitecturas resilientes \emph{offline-first}.
    \item \textbf{Valores Corporativos:} Rigor metodológico, seguridad por diseño (\emph{Zero Trust}), transparencia auditable y compromiso irrestricto con la continuidad de servicio de nuestros mandantes.
    \item \textbf{Presencia Geográfica, Instalaciones y Cobertura de Soporte en Terreno:} Casa matriz y centro de ingeniería de software ubicado en Viña del Mar (Región de Valparaíso), y oficina de operaciones técnicas, pruebas de banco y laboratorio de dispositivos IoT/embarcados en Santiago (Región Metropolitana), lo que asegura una capacidad de despliegue y soporte presencial expedito sobre el eje logístico central Valparaíso -- Santiago -- San Bernardo -- Los Andes. Adicionalmente, para dar estricto cumplimiento a las exigencias del Capítulo 8.4 (reemplazo de hardware crítico en terreno en un máximo de cuatro horas) y al requerimiento RT-21.16 (atención técnica y movilización inmediata a sitios alejados) de las Bases Técnicas Transversales \parencite{bases_transversales_2026}, audIT mantiene convenios de servicio vigentes con una red técnica certificada con presencia permanente en los cuatro nodos regionales de Transportes Curimón S.A. (\textbf{Antofagasta, Talca, Los Ángeles y Puerto Montt}), garantizando capacidad de intervención física in situ y reemplazo de hardware telemático con tiempos de respuesta inferiores a cuatro (4) horas a lo largo de todo el corredor de 3.000 kilómetros de la Ruta 5.
\end{itemize}

Las capacidades instaladas de la compañía se estructuran en tres líneas de negocio y un catálogo de productos y servicios especializados de misión crítica:

\begin{enumerate}
    \item \textbf{Ingeniería IoT y Sistemas de Terreno (Edge Computing):} Desarrollo de firmware industrial, integración de telemetría vehicular multi-marca (SAE J1939, rFMS y tacógrafos digitales) y provisión del paquete de software embarcado \emph{audIT EdgeHub} con resiliencia, almacenamiento local inalterable y sincronización automática tras desconexión prolongada.
    \item \textbf{Plataformas de Datos y Analítica Avanzada:} Modelamiento e implantación de repositorios operacionales y analíticos en nube híbrida (\emph{audIT TelemetryCore}), ingestión masiva de eventos telemáticos en tiempo real, motores de despacho algorítmico, reportería analítica y tableros de control operacional en tiempo real.
    \item \textbf{Consultoría en Arquitectura Híbrida y Ciberseguridad Operacional:} Diseño de capas anticorrupción (\emph{Domain-Driven Design}) para convivencia con sistemas ERP heredados, auditorías de seguridad perimetral, consultoría de desarrollo seguro (OWASP SAMM) y gobierno de datos personales bajo la Ley N.\textordmasculine{} 21.719 \parencite{ley21719}.
\end{enumerate}

\subsection{Estructura organizacional, dotación, certificaciones y alianzas}

audIT cuenta con una dotación permanente de \textbf{22 profesionales de planta} (ingenieros civiles informáticos, ingenieros de software, ingenieros de telecomunicaciones/electrónica y científicos de datos), complementada según demanda por una red certificada de técnicos de instalación en terreno. La organización adopta un modelo de gestión matricial y de ingeniería colaborativa, estructurado de manera proporcionada y altamente especializada:

\figuraAncha{\subdocRuta/Img.png}{Estructura organizacional y dotación de audIT}{organigrama}

La gobernanza corporativa y técnica de audIT reside en \textbf{3 Roles de Dirección Ejecutiva}, quienes concentran el liderazgo metodológico y la toma de decisiones críticas:

\begin{itemize}
    \item \textbf{Gerencia General / Dirección de Proyectos (Asignado al Rol Habilitante: Jefe de Proyecto):} Ingeniero Civil Informático, certificado PMP\textsuperscript{\textregistered} (\emph{Project Management Professional}), con más de 12 años liderando contratos de modernización tecnológica e integración de sistemas críticos en transporte e infraestructura logística.
    \item \textbf{Dirección de Arquitectura y Datos (Asignado a los Roles Habilitantes: Arquitecto de Solución y Líder de Datos):} Máster en Ciencias de la Computación, certificado \emph{Microsoft Certified: Azure Solutions Architect Expert} y \emph{CDMP (Certified Data Management Professional)}, responsable de la arquitectura distribuida, capas de interoperabilidad y gobierno del dato.
    \item \textbf{Jefatura de IoT y Terreno (Asignado al Rol Habilitante: Líder de Operación):} Ingeniero en Telecomunicaciones y Electrónica, especialista en integración telemática embarcada bajo estándares SAE J1939, rFMS y protocolos industriales de campo.
\end{itemize}

La dotación operativa se compone de \textbf{19 profesionales de planta}, distribuidos en tres unidades clave donde se asignan los roles habilitantes restantes exigidos por el Artículo 34.1 de las Bases Administrativas \parencite{bases_admin_2026}:

\begin{itemize}
    \item \textbf{12 Ingenieros de Software, Firmware y Datos:} Organizados en células y duplas operacionales pareadas (\emph{pair-engineering}), asegurando continuidad operacional, cobertura cruzada y ausencia de dependencia de individuos únicos tanto en backend, servicios en la nube y analítica, como en firmware embarcado y telemetría de terreno.
    \item \textbf{4 Especialistas en Aseguramiento de Calidad, Ciberseguridad y Auditoría Interna ISO (Asignados a los Roles Habilitantes: Líder de Calidad y Encargado de Seguridad de la Información):} Dedicados de forma transversal a los controles de calidad (\emph{quality gates}), matrices de trazabilidad, pruebas de intrusión y homologación en terreno, y auditoría continua de procesos bajo estándares internacionales y la Ley N.\textordmasculine{} 21.719 \parencite{ley21719}.
    \item \textbf{3 Analistas de Gestión de Proyectos (PMO) y Soporte:} A cargo de la planificación operativa, seguimiento de hitos, control contractual, compras técnicas y mesa de soporte especializado de segundo nivel 24/7.
\end{itemize}

\textbf{Certificaciones Individuales del Personal de Planta (Bases Técnicas FEP02.26 · Cap. 15.3 y RT-15.07 a RT-15.09):}
audIT garantiza que su dotación permanente de 22 profesionales cubre con holgura las certificaciones individuales mínimas obligatorias exigidas para el Contrato, manteniéndolas plenamente vigentes:
\begin{itemize}
    \item \emph{Gestión de Proyectos (PMP / PRINCE2):} 2 profesionales certificados (Gerente General y Analista PMO Senior).
    \item \emph{Gestión de Servicios TI (ITIL 4 Foundation o superior):} 5 profesionales certificados (3 Analistas PMO/Soporte y 2 Especialistas QA).
    \item \emph{Gobierno de TI (COBIT o equivalente):} 2 profesionales certificados (Director de Proyectos y Especialista de Auditoría ISO).
    \item \emph{Arquitectura de Nube del Proveedor Ofertado (Azure Solutions Architect Expert):} 3 ingenieros certificados (Director de Arquitectura y 2 Ingenieros Senior de Célula Cloud).
    \item \emph{Seguridad de la Información (CISSP / CISM / CEH):} 2 profesionales certificados (Especialistas de Ciberseguridad del área de Calidad y Seguridad).
    \item \emph{Bases de Datos del Motor Ofertado (PostgreSQL / Azure Data Engineer):} 2 ingenieros certificados (Líder de Datos e Ingeniero Senior Backend).
    \item \emph{Calidad y Pruebas de Software (ISTQB Certified Tester):} 2 especialistas certificados (Líder de Calidad y Especialista QA).
\end{itemize}

\textbf{Certificaciones Institucionales, Alianzas Tecnológicas y Competencias Sectoriales:}

\begin{itemize}
    \item \textbf{ISO 9001:2015 (Sistema de Gestión de la Calidad):} Certificación institucional plenamente vigente para el diseño, desarrollo, implantación y soporte de soluciones de software e integración de hardware de terreno \parencite{iso9001}.
    \item \textbf{ISO 27001:2022 (Sistema de Gestión de Seguridad de la Información):} En fase final de certificación (Fase 2 superada exitosamente y emisión de certificado en curso), debidamente respaldada bajo el Plan de Certificación Institucional con hitos verificables conforme a lo dispuesto en el Artículo 34.1 de las Bases Administrativas \parencite{bases_admin_2026}, con alcance sobre el tratamiento, almacenamiento seguro y transmisión de datos telemáticos, información sensible de mandantes y trazabilidad contractual \parencite{iso27001}.
    \item \textbf{Competencias Sectoriales Acreditadas del Adjudicatario (FEP03.10.26 · RT-15.02):} audIT y su equipo técnico cuentan con conocimiento acreditado del régimen laboral especial de jornada del conductor de carga terrestre (Código del Trabajo, Artículo 25 bis \parencite{codigo_trabajo_25bis}), de la normativa reglamentaria para el transporte y almacenamiento de sustancias peligrosas (Decreto Supremo N.\textordmasculine{} 298 del Ministerio de Transportes y Telecomunicaciones \parencite{ds298} y Decreto Supremo N.\textordmasculine{} 43 del Ministerio de Salud \parencite{ds43}), y experiencia empírica demostrada en telemetría de flotas pesadas con tolerancia certificada a operación desconectada prolongada.
    \item \textbf{Certificaciones Institucionales Complementarias y Deseables (FEP02.26 · Cap. 15.2):} La compañía opera bajo procesos formalmente alineados con los estándares ISO/IEC 20000-1 \parencite{iso20000_1} (Gestión de Servicios TI) e ISO 22301 \parencite{iso22301} (Continuidad del Negocio), garantizando la resiliencia operativa en contratos de misión crítica.
    \item \textbf{Alianzas Estratégicas y Capacidad de Proveedor Cloud:} \emph{Cloud Solution Provider (CSP) Tier 1 (Direct Partner)} de \textbf{Microsoft Azure} con especialización avanzada en infraestructura híbrida (\emph{Azure Arc}, \emph{Azure IoT Hub} y bases de datos distribuidas) y convenios de homologación e integración directa con fabricantes de hardware telemático industrial de estándar automotriz.
\end{itemize}

\subsection{1.2.1 Capacidades Técnicas Instaladas, Stack Tecnológico y Metodologías de Trabajo}

audIT domina un ecosistema tecnológico moderno, robusto y probado en escenarios de alta exigencia operacional:

\begin{itemize}
    \item \textbf{Nivel Embarcado y Terreno:} Desarrollo en C/C++, Rust y Python embebido sobre sistemas Linux industrial (\emph{Yocto/Debian Embedded}); almacenamiento local indexado mediante bases de datos transaccionales ultraligeras (\emph{SQLite en modo WAL}) con búfer circular inalterable. Estas capacidades sustentan tanto el software de pasarelas vehiculares como la interoperabilidad con aplicaciones móviles de cabina multiplataforma.
    \item \textbf{Nivel de Integración y Servicios:} Arquitecturas basadas en microservicios desacoplados vía \emph{APIs RESTful} seguras (\textbf{OpenAPI 3.1}), \emph{brokers} de mensajería orientados a eventos (\emph{Apache Kafka / Azure Event Hubs / Azure Service Bus}) y capas anticorrupción (\emph{Domain-Driven Design}) para aislar componentes legados y ERPs existentes.
    \item \textbf{Nivel de Nube y Datos:} Contenerización con \emph{Docker / Kubernetes (AKS)}, bases de datos híbridas relacionales y temporales (\emph{PostgreSQL / TimescaleDB}), almacenamiento estructurado inmutable (\emph{Azure Blob Storage WORM}), y tableros analíticos en tiempo real (\emph{Power BI Embedded} y servicios analíticos dedicados en arquitectura Lakehouse).
\end{itemize}

\textbf{Metodologías de Trabajo Certificadas:}

\begin{itemize}
    \item \emph{Desarrollo de Software y Analítica:} Enfoque ágil \textbf{Scrum / Kanban} con integración y despliegue continuo (\textbf{CI/CD}), asegurando entregas funcionales iterativas validadas por el usuario.
    \item \emph{Componentes Críticos e Integración de Hardware:} Metodología ingenieril en cascada adaptada (\textbf{Modelo V} conforme a ISO/IEC 12207 \parencite{iso12207}) para el aseguramiento de firmware, homologación vehicular y compatibilidad eléctrica en cabina.
\end{itemize}

\section{1.3 Experiencia Relevante en la Industria del Caso y en Proyectos de Complejidad Equivalente}

La experiencia de audIT aborda de forma directa y verificable las complejidades del sector logístico y del transporte de carga terrestre:

\begin{enumerate}
    \item \textbf{Telemetría Vehicular y Geocercas en Tiempo Real:} En el proyecto ejecutado para \emph{Transportes del Sur Ltda.}, audIT diseñó e implementó un sistema híbrido de telemetría a bordo sobre una flota de 120 tracto-camiones de carga interurbana, integrando posicionamiento en tiempo real, automatización de la captura de eventos de conducción, control de geocercas en terminales y despacho centralizado. La arquitectura combinó \emph{edge computing} a bordo del vehículo con procesamiento analítico en nube Azure, alcanzando un SLA de disponibilidad del 99,5\,\% mensual y procesando aproximadamente 5.500 eventos telemáticos diarios.
    
    \item \textbf{Operación Desconectada y Sincronización Resiliente:} En \emph{AgroFrutícola Los Andes S.A.}, audIT desplegó una arquitectura \emph{offline-first} sobre 6 faenas cordilleranas con intermitencia severa de red móvil, operando 1.500 sensores telemáticos activos con almacenamiento local y reconciliación transaccional determinista sin pérdida de datos tras períodos de desconexión de hasta 72 horas. La arquitectura híbrida (nodos concentradores locales on-premise + nube) alcanzó un SLA de disponibilidad del 99,0\,\% con tolerancia certificada a operación desconectada prolongada.
    
    \item \textbf{Escalabilidad Híbrida y Alta Disponibilidad en Logística:} En el contrato de modernización de \emph{Logística y Distribución Multimodal Bicentenario S.A.}, audIT diseñó una plataforma analítica híbrida (nube pública Microsoft Azure + servidores locales en plantas de transferencia) con SLA certificado del 99,6\,\% mensual, procesando aproximadamente 2,5 TB de transacciones logísticas mensuales con más de 800 usuarios concurrentes, incluyendo ingestión de eventos de despacho, reportería analítica ejecutiva y auditoría transaccional en tiempo real.
\end{enumerate}

En conjunto, estos tres proyectos evidencian que audIT no solo conoce la industria del transporte de carga, sino que ha resuelto previamente los tres problemas técnicos centrales que caracterizan los desafíos del sector: telemetría vehicular con visibilidad de flota en tiempo real, operación confiable en escenarios de conectividad intermitente o inexistente, y arquitectura de datos escalable en nube híbrida con niveles de servicio exigentes.

\begin{tablaAudit}{Formulario T-6 --- Experiencia en Proyectos Similares}{t6_experiencia}{@{}p{2.3cm}>{\raggedright\arraybackslash}p{4.1cm}>{\raggedright\arraybackslash}p{4.1cm}>{\raggedright\arraybackslash}p{4.1cm}@{}}
\textbf{Campo Reglamentario} & \textbf{Proyecto 1 (Experiencia Específica Híbrida y Transporte)} & \textbf{Proyecto 2 (Resiliencia Offline-First y Terreno)} & \textbf{Proyecto 3 (Experiencia Específica SLA $\ge 99{,}5\%$ y Datos)} \\ \midrule
\textbf{1. Nombre del proyecto} & Sistema Híbrido de Telemetría, Trazabilidad y Control de Flota & Plataforma IoT y Telemetría Resiliente con Arquitectura \emph{Offline-First} & Plataforma Analítica Centralizada y Gestión Logística en Alta Disponibilidad \\ \midrule
\textbf{2. Cliente / mandante} & Transportes del Sur Ltda. & AgroFrutícola Los Andes S.A. & Logística y Distribución Multimodal Bicentenario S.A. \\ \midrule
\textbf{3. Industria} & Transporte terrestre de carga interurbana & Agroindustria, faenas agrícolas y cadena de frío & Logística, distribución multimodal y servicios de carga \\ \midrule
\textbf{4. Año de inicio y de término} & 2023 -- 2024 (14 meses, finalizado y en operación) & 2022 -- 2023 (12 meses, finalizado y en operación) & 2024 -- 2025 (15 meses, finalizado y en operación) \\ \midrule
\textbf{5. Monto del contrato (rango)*} & Rango: UF 3.000 -- UF 5.000 & Rango: UF 2.000 -- UF 3.500 & Rango: UF 4.500 -- UF 6.500 \\ \midrule
\textbf{6. Alcance ejecutado} & Diseño, provisión e implantación de telemetría a bordo (120 camiones), control de eventos de conducción, geocercas en terminales y despacho centralizado. & Despliegue de concentradores IoT y telemetría en 6 faenas remotas sin señal móvil; almacenamiento local y reconciliación sin pérdida tras reconexión. & Diseño de plataforma en nube híbrida (Azure + nodos locales), ingestión de eventos de despacho, reportería analítica ejecutiva y auditoría transaccional. \\ \midrule
\textbf{7. Arquitectura (nube / on-premise / híbrida)} & \textbf{Híbrida} (Edge computing a bordo en camión + nube Azure para analítica) & \textbf{Híbrida} (Nodos concentradores locales on-premise + nube) & \textbf{Híbrida} (Nube pública Microsoft Azure + servidores locales de enlace en plantas de transferencia) \\ \midrule
\textbf{8. Nivel de servicio comprometido} & Disponibilidad de plataforma: \textbf{99,5\,\%} mensual (24/7 contractualmente). & Disponibilidad: \textbf{99,0\,\%}; tolerancia a operación desconectada de \textbf{72 horas} sin pérdida de datos. & Disponibilidad de plataforma: \textbf{99,6\,\%} mensual garantizada contractualmente. \\ \midrule
\textbf{9. Volumen de operación soportado} & 120 tracto-camiones; $\approx 5.500$ eventos telemáticos diarios procesados. & 6 faenas remotas; 1.500 sensores telemáticos activos; búfer de millones de registros locales. & $\approx 2{,}5\text{ TB}$ mensuales de transacciones logísticas; $> 800$ usuarios concurrentes. \\ \midrule
\textbf{10. Rol de la empresa (principal / socio)} & Principal (Contratista Principal, 100\,\% de la ingeniería y desarrollo) & Principal (Contratista Principal, 100\,\% de la ingeniería y despliegue) & Principal (Contratista Principal, 100\,\% de la ingeniería y desarrollo) \\ \midrule
\textbf{11. Contraparte de referencia y contacto} & Marcelo Iturra V.\newline\textit{Gerente de Operaciones}\newline Correo: {\footnotesize\url{miturra@transportesdelsur.cl}}\newline Fono: +56 9 7845 1290 & Paula Concha M.\newline\textit{Jefa de TI}\newline Correo: {\footnotesize\url{pconcha@agrofruticolalosandes.cl}}\newline Fono: +56 9 8451 9023 & Rodrigo Baeza S.\newline\textit{Subgerente Corp. de Sistemas}\newline Correo: {\footnotesize\url{rbaeza@logisticabicentenario.cl}}\newline Fono: +56 9 6521 3487 \\
\end{tablaAudit}

\emph{*Nota de Blindaje Económico (Bases Administrativas FEP01.26 $\cdot$ Artículos 50.2$^\circ$ y 53$^\circ$): Los montos indicados corresponden estrictamente a rangos de contratos pretéritos ejecutados y liquidados con terceros en los períodos señalados. En estricto cumplimiento del Artículo 50.2$^\circ$ de las Bases Administrativas \parencite{bases_admin_2026}, no constituyen tarifas, honorarios ni precios de la oferta técnica ni económica del presente proceso licitatorio.}

\section{1.4 Modelo de Gobierno Interno de Calidad, Seguridad y Gestión del Conocimiento}

El gobierno interno de audIT se ejecuta a través del área transversal de Calidad y Seguridad de la Información, y se sostiene en tres pilares operativos:

\begin{enumerate}
    \item \textbf{Aseguramiento de Calidad (ISO 9001):} Cada fase de desarrollo cuenta con hitos de compuerta (\emph{quality gates}). Ningún entregable pasa a producción sin la aprobación formal del Responsable de Calidad (rol independiente de los líderes de desarrollo), quien audita el cumplimiento de matrices de trazabilidad de requisitos, cobertura de pruebas unitarias y de integración ($\ge 85\%$) y actas formales de homologación en terreno.
    
    \item \textbf{Seguridad y Privacidad por Diseño (ISO 27001, FEP02 Cap. 11.5 y Leyes N.\textordmasculine{} 21.719 y 21.663):} Aplicación estricta de Arquitectura de Confianza Cero conforme a \textbf{NIST SP 800-207} \parencite{nist_sp800_207} y el marco \textbf{NIST Cybersecurity Framework (CSF) 2.0} \parencite{nist_csf2}. La seguridad aplicativa se gobierna bajo \textbf{OWASP ASVS 4.0 nivel 2} \parencite{owasp_asvs}, OWASP Top 10 y endurecimiento de sistemas mediante \textbf{CIS Benchmarks} \parencite{cis_benchmarks}, asegurando la cadena de suministro de software bajo \textbf{SLSA nivel 3} (RT-11.28). La información telemática y transaccional se cifra de extremo a extremo: AES-256 en reposo y TLS 1.3 en tránsito. Conforme al requisito \textbf{RT-11.10}, se aplica \textbf{cifrado a nivel de campo (\emph{field-level encryption})} sobre los datos personales de los 258 conductores externos, registros de jornada y las condiciones tarifarias pactadas con los 148 transportistas subcontratados. El registro de auditoría es criptográficamente inalterable (RT-16.07), registrando todo acceso a información sensible (RT-16.09), garantizando plena oponibilidad jurídica bajo la Ley N.\textordmasculine{} 21.719 \parencite{ley21719} de Protección de Datos Personales y la Ley N.\textordmasculine{} 21.663 \parencite{ley21663} (Ley Marco de Ciberseguridad).
    
    \item \textbf{Gestión del Conocimiento y Mitigación de Dependencia:} audIT mantiene una base de conocimiento institucional automatizada (arquitecturas de referencia, patrones de integración de buses vehiculares, guías de resolución de fallas y \emph{post-mortems} de incidentes). La estandarización de artefactos y el trabajo en parejas de ingeniería ---organizadas en células y duplas operacionales--- evitan silos cognitivos o dependencia de individuos específicos, garantizando soporte ininterrumpido a nuestros clientes aun ante rotación de personal.
\end{enumerate}
